\chapter{Conclusão}
\label{cap:conclusao}

Este trabalho descreveu o projeto, o desenvolvimento e a validação técnica do
Laboratório 404, um jogo sério 3D no qual uma inteligência artificial local
participa do mundo por meio de diálogo, sem controlá-lo. O objetivo geral ---
desenvolver e validar tecnicamente um jogo dessa natureza, executável em
computador pessoal sem GPU dedicada e com apoio procedural ao ensino de
automação industrial --- foi atingido no nível de protótipo: as seis etapas
previstas foram concluídas e demonstráveis; a suíte automatizada cobre 211
verificações do motor mais 16 testes do adaptador, todas aprovadas; o
desempenho atende ao critério de hardware modesto; e a integração com a IA
funciona de ponta a ponta com degradação segura verificada.

Os objetivos específicos foram cumpridos integralmente: a especificação com
requisitos e critérios de aceitação foi escrita e mantida como contrato
(\autoref{ap:requisitos}); o desenvolvimento incremental em provas de
conceito foi executado com evidências a cada etapa (\autoref{cap:desenvolvimento});
o pipeline de assets do Blender ao motor foi implementado com convenções de
escala, nomenclatura e colisão; a IA local foi integrada em processo separado,
com contrato validado, intenções restritas e \emph{fallback}; o mundo reativo
e a geração procedural foram implementados e testados; a suíte de validação e
o arnês de evidências foram construídos e executados; e as medições de
desempenho e latência foram realizadas com limitações declaradas
(\autoref{cap:resultados}).

As contribuições principais são o padrão de \emph{IA diegética com capacidade
limitada}, o processo incremental rastreável adequado a desenvolvimento
individual assistido por agentes de IA, o artefato aberto e documentado e o
registro de lições de engenharia e de ambiente que afetam a reprodutibilidade
de projetos semelhantes (\autoref{sec:licoes}).

As limitações permanecem claras e declaradas: ausência de estudo com usuários,
amostra pequena de latências, modelo-alvo não executado, validações manuais de
hardware pendentes, empacotamento independente não gerado e edições
normativas ainda a conferir. Nenhuma delas invalida os resultados
apresentados, mas todas delimitam o alcance das conclusões: o que este
trabalho demonstra é viabilidade técnica e rastreabilidade de processo, não
eficácia educacional.

Como desdobramento imediato, foram propostos seis trabalhos futuros, do
benchmark do modelo-alvo ao estudo controlado com estudantes
(\autoref{sec:futuros}). O caminho natural é que a próxima etapa deste
trabalho não seja construir mais jogo, e sim buscar evidência de aprendizagem:
um instrutor ficcional que responde em segundos, não controla nada e está
sempre disponível já existe; falta medir o que ele ensina.
